\chapter*{\centering Словарь терминов}
\addcontentsline{toc}{chapter}{Словарь терминов}


\textbf{Никелевые жаропрочные сплавы} -
сложнолегированные металлические материалы на основе никеля, предназначенные для работы при высоких температурах, длительных нагрузках и в агрессивных средах. Характеризуются высокой жаропрочностью, коррозионной стойкостью и структурной стабильностью.

\textbf{Суперсплавы} -
класс высоколегированных сплавов, сохраняющих механические свойства при температурах, близких к температуре плавления. В данной работе под суперсплавами преимущественно понимаются никелевые жаропрочные сплавы.

\textbf{$\gamma$-матрица} -
основная фаза никелевых жаропрочных сплавов, представляющая собой твердый раствор на основе никеля с гранецентрированной кубической решеткой. В $\gamma$-матрице распределены упрочняющие фазы.

\textbf{$\gamma'$-фаза} -
упрочняющая интерметаллидная фаза, она является одним из основных факторов высокой жаропрочности никелевых суперсплавов.

% \textbf{$\gamma''$-фаза} -
% упрочняющая метастабильная интерметаллидная фаза, характерная для некоторых никелевых сплавов, содержащих ниобий. Вносит вклад в повышение прочности, особенно при промежуточных температурах.

\textbf{Температура сольвуса} -
температура, выше которой определенная фаза полностью растворяется в матрице сплава. В данной работе особое значение имеет температура сольвуса $\gamma'$-фазы, обозначаемая как $T_{\mathrm{solvus}}$.

\textbf{$T_{\mathrm{solvus}}$} -
обозначение температуры сольвуса. Используется как целевая переменная при моделировании фазовой стабильности никелевых суперсплавов.

% \textbf{Фазовая стабильность} -
% способность фазовой структуры сплава сохраняться при высокотемпературном воздействии. Чем выше фазовая стабильность, тем медленнее происходит растворение, укрупнение или деградация упрочняющих фаз.

% \textbf{Жаропрочность} -
% способность материала сохранять прочность и сопротивляться разрушению при длительном воздействии высоких температур и механических нагрузок.

% \textbf{Термическая стабильность} -
% способность структуры и свойств материала сохраняться при воздействии повышенной температуры в течение длительного времени.

\textbf{Легирование} -
введение в сплав дополнительных химических элементов с целью изменения его структуры и свойств. В никелевых суперсплавах легирование используется для повышения прочности, жаростойкости, коррозионной стойкости и фазовой стабильности.

% \textbf{Легирующий элемент} -
% химический элемент, добавляемый в основной металл для улучшения свойств сплава. В никелевых суперсплавах такими элементами являются, например, Al, Ti, Cr, Co, Mo, W, Ta, Re, Nb, Hf, C, B и др.

% \textbf{Балансный никель} -
% никель, содержание которого определяется как остаток до 100\% после учета всех легирующих элементов. В моделях никель часто не используется как отдельный входной признак, поскольку является основой сплава.

% \textbf{Массовая доля} -
% содержание элемента в сплаве, выраженное в процентах от общей массы. Обычно обозначается как wt.\%.

% \textbf{Атомная доля} -
% содержание элемента, выраженное как отношение числа атомов данного элемента к общему числу атомов в сплаве. Более физически корректна для описания фазообразования и стехиометрии.

% \textbf{Нормализация данных} -
% преобразование численных признаков к удобному диапазону значений. В данной работе используется нормализация концентраций легирующих элементов относительно никеля.

% \textbf{Нормализация относительно никеля} -
% способ представления состава, при котором атомная доля каждого легирующего элемента делится на атомную долю никеля:
% $$     z_i = \frac{x_i}{x_{\mathrm{Ni}}}.     $$
% Такое преобразование позволяет исключить никель из входного вектора модели и привести признаки к сопоставимому масштабу.

\textbf{Параметр Ларсона--Миллера} -
температурно-временной параметр, объединяющий температуру и время выдержки или разрушения:
$$     P_{LM}=T(C+\log_{10}\tau).     $$
Используется для описания длительной прочности и деградации материала при высоких температурах.

% \textbf{$P_{LM}$}
% обозначение параметра Ларсона--Миллера. В отчете используется для описания зависимости прочности от условий испытания.

% \textbf{$P_{LM}^{*}$}
% нормированная форма параметра Ларсона--Миллера:
% $$     P_{LM}^{*}=T(20+\log_{10}\tau)\cdot 10^{-5}.     $$
% Нормировка используется для согласования масштаба данного параметра с концентрациями элементов во входном векторе нейронной сети.

\textbf{Предел прочности} -
максимальное механическое напряжение, которое материал выдерживает до разрушения. Обозначается как $\sigma$.

\textbf{$\sigma$} -
обозначение прочности или предела прочности сплава. В одной из моделей является целевой переменной, предсказываемой нейронной сетью.

\textbf{$\sigma(P_{LM})$}
зависимость прочности от параметра Ларсона--Миллера. Отражает изменение прочности сплава при различных температурно-временных условиях.

% \textbf{Логарифмическое преобразование} -
% преобразование целевой переменной, применяемое для уменьшения диапазона значений и повышения устойчивости обучения модели. Для прочности используется выражение:
% $$     y=-\log_{10}\sigma.     $$

% \textbf{Искусственная нейронная сеть} -
% вычислительная модель, состоящая из связанных между собой нейронов и способная аппроксимировать сложные нелинейные зависимости между входными и выходными данными.

% \textbf{ANN} -
% Artificial Neural Network --- искусственная нейронная сеть.

% \textbf{Глубокая нейронная сеть} -
% нейронная сеть, содержащая несколько скрытых слоев. Позволяет извлекать более сложные иерархические зависимости из данных.

% \textbf{Многослойный персептрон} -
% тип нейронной сети прямого распространения, состоящий из входного слоя, одного или нескольких скрытых слоев и выходного слоя.

% \textbf{Скрытый слой} -
% внутренний слой нейронной сети, в котором выполняется нелинейное преобразование входных данных.

% \textbf{Выходной слой} -
% последний слой нейронной сети, формирующий итоговое предсказание модели.

% \textbf{Функция активации} -
% нелинейная функция, применяемая к выходу нейрона. Позволяет нейронной сети моделировать нелинейные зависимости.

% \textbf{Обучение нейронной сети} -
% процесс подбора весов и смещений сети с целью минимизации ошибки между предсказанными и экспериментальными значениями.

% \textbf{Обратное распространение ошибки} -
% алгоритм вычисления градиентов ошибки по параметрам нейронной сети, используемый при ее обучении.

% \textbf{Переобучение} -
% ситуация, при которой модель слишком хорошо запоминает обучающую выборку, но плохо предсказывает новые данные.

% \textbf{Регуляризация} -
% метод ограничения сложности модели для уменьшения переобучения и повышения обобщающей способности.

% \textbf{Байесовская регуляризация} -
% метод регуляризации, при котором обучение нейронной сети рассматривается с вероятностной точки зрения, а в функцию ошибки добавляется штраф за большие значения весов.

\textbf{BRANN} -
Bayesian Regularized Artificial Neural Network --- искусственная нейронная сеть с байесовской регуляризацией. Используется для повышения устойчивости модели к переобучению.

% \textbf{Ансамбль нейронных сетей} -
% совокупность нескольких независимо обученных нейронных сетей, прогнозы которых объединяются, например, усреднением. Позволяет повысить стабильность и точность предсказаний.

% \textbf{Бутстреп} -
% метод многократного случайного формирования обучающих и проверочных выборок из исходных данных. Используется для повышения надежности оценки модели.

\textbf{Обучающая выборка} -
часть данных, используемая для настройки параметров модели.

\textbf{Валидационная выборка} -
часть данных, используемая для контроля качества модели во время обучения и предотвращения переобучения.

\textbf{Тестовая выборка} -
независимая часть данных, используемая для окончательной оценки качества обученной модели.

\textbf{Априорная информация} -
заранее известные сведения о физической природе объекта моделирования. В данной работе это информация о роли легирующих элементов в структуре и свойствах суперсплавов.

\textbf{Матрица связей} -
матрица, описывающая принадлежность легирующих элементов к различным группам влияния. Используется для внедрения физической информации в архитектуру нейронной сети.

\textbf{Дифференциальный слой} -
слой нейронной сети, передающий информацию о чувствительности выходного параметра к изменению входных признаков. Применяется для повышения интерпретируемости модели.

\textbf{MSE} -
Mean Squared Error --- среднеквадратичная ошибка:
$$     MSE =     \frac{1}{n}     \sum_{i=1}^{n}     \left(y_i-\hat{y}_i\right)^2.     $$

\textbf{RMSE} -
Root Mean Squared Error --- корень из среднеквадратичной ошибки. В отчете также используется относительная форма RMSE для оценки ошибки прогноза прочности.

\textbf{RMAE} -
Relative Mean Absolute Error --- относительная средняя абсолютная ошибка. Используется для оценки точности прогноза температуры сольвуса.

\textbf{RRMSE} -
Relative Root Mean Squared Error --- относительная среднеквадратичная ошибка. Используется для оценки качества модели при прогнозировании температуры сольвуса.

\textbf{Химический состав} -
совокупность концентраций элементов, входящих в сплав. В моделях машинного обучения химический состав используется как входной вектор признаков.

\textbf{Вектор признаков} -
набор численных характеристик объекта, передаваемый на вход модели. Для суперсплавов включает концентрации легирующих элементов и, при необходимости, параметр Ларсона--Миллера.

\textbf{Целевая переменная} -
величина, которую необходимо предсказать с помощью модели. В данном отчете целевыми переменными являются температура сольвуса $T_{\mathrm{solvus}}$ и прочность $\sigma$.